\documentclass{article}

\usepackage{amsmath}
%\usepackage{amsfonts}
\usepackage{amsthm}
%\usepackage{amssymb}
%\usepackage{mathrsfs}
%\usepackage{fullpage}
%\usepackage{mathptmx}
%\usepackage[varg]{txfonts}
\usepackage{color}
\usepackage[charter]{mathdesign}
\usepackage[pdftex]{graphicx}
%\usepackage{float}
%\usepackage{hyperref}
%\usepackage[modulo, displaymath, mathlines]{lineno}
%\usepackage{setspace}
%\usepackage[titletoc,toc,title]{appendix}
\usepackage{natbib}
\usepackage[makeroom]{cancel}
\usepackage[only,llbracket,rrbracket]{stmaryrd}

%\linenumbers
%\doublespacing

\theoremstyle{definition}
\newtheorem*{defn}{Definition}
\newtheorem*{exm}{Example}

\theoremstyle{plain}
\newtheorem*{thm}{Theorem}
\newtheorem*{lem}{Lemma}
\newtheorem*{prop}{Proposition}
\newtheorem*{cor}{Corollary}

\newcommand{\argmin}{\text{argmin}}
\newcommand{\ud}{\hspace{2pt}\mathrm{d}}
\newcommand{\bs}{\boldsymbol}
\newcommand{\PP}{\mathsf{P}}
\let\divsymb=\div % rename builtin command \div to \divsymb
\renewcommand{\div}[1]{\operatorname{div} #1} % for divergence
\newcommand{\Id}[1]{\operatorname{Id} #1}

\title{Simulating the flow of viscous, incompressible fluids with a free surface using terrain-following coordinates}
\author{Daniel R. Shapero, Emily Wilbur}
\date{}

\begin{document}

\maketitle

% --------------------
\section{Introduction}

Here we describe discretization schemes for viscous free-surface fluid flows.
In atmospheric science and oceanography, using terrain-following coordinates is a common approach for handling variable surface and bottom topography.
For these problems, viscous dissipation is negligible to a good approximation.
The equations of motion contain only first-order derivatives in space and time.
Here we are interested in other free-surface fluid flows in the earth sciences, such as the flow of glaciers or lava.
Viscous effects are dominant for these types of flows and so the equations of motion contain second-order derivatives.
This makes it much harder to apply terrain-following coordinates because the metric terms are discontinuous and live inside a derivative.
Here we show how to work around this difficulty using discontinuous Galerkin methods.

Using interior penalty discontinuous Galerkin methods requires setting certain penalty parameters.
Many papers in the research literature assume that this constant is big enough without deriving how big that is.
An essential tool for us will be the explicit estimates for constants in the finite element trace inverse inequality of \citet{warburton2003constants}.
This will let us set the constants explicitly in terms of the mesh regularity.

We assume that the fluid flow takes place over a variable topography with elevation $b$.
The fluid surface is at height $s$ and the film thickness is $h = s - b$.

% -----------------------------------
\section{DG discretization of Stokes}

The incompressible Stokes equations can be derived through minimization of the free energy dissipation rate
\begin{equation}
    G(u, p) = \int_\Omega\left\{\frac{1}{2}\tau :\dot\varepsilon - p\nabla\cdot u - f\cdot u\right\}\mathrm dx
\end{equation}
where $\dot\varepsilon = \frac{1}{2}(\nabla u + \nabla u^*)$ is the rate of strain tensor and $\tau = 2\mu\dot\varepsilon$ is the deviatoric stress tensor.
We will use the minimization form in the following because it requires fewer characters to type and because I like it.

The description above is not quite complete because we also need to specify boundary conditions.
We will concern ourselves with two types of boundary conditions.
First, we can prescribe the boundary traction by adding a term to the free energy rate:
\begin{equation}
    \ldots - \int_\Gamma \tau_\Gamma : n\otimes u\,\mathrm d\gamma.
\end{equation}
We might instead want to fix the value of the velocity.
The conventional approach is to modify the linear system at the time of solution to enforce this Dirichlet condition.

Nitsche's method is a way to instead enforce Dirichlet conditions by modifying the variational form.
One way of looking at Nitsche's method is as an augmented Lagrangian scheme together with an exact solution for the Lagrange multipliers.
The Nitsche functional is
\begin{equation}
    G_\alpha(u, p) = G(u, p) - \int_\Gamma(\tau - pI) : n\otimes(u - u_\Gamma)\mathrm d\gamma + \int_\Gamma\frac{\alpha\mu}{2\ell}|u - u_\Gamma|^2\mathrm d\gamma.
\end{equation}
Here $\ell$ is a measure of the local facet diameter and $\alpha$ is a constant related to the mesh regularity.

The symmetric interior penalty discontinuous Galerkin method is what happens when we apply Nitsche's method to every cell $\omega$ of the mesh:
\begin{align}
    & G_\alpha(u, p) = \sum_\omega\int_\omega\left\{\frac{1}{2}\tau:\dot\varepsilon - p\nabla\cdot u - f\cdot u\right\}\mathrm dx \nonumber \\
    & \qquad - \sum_\gamma\int_\gamma \left\langle 2\mu\dot\varepsilon - pI \right\rangle: [u\otimes n]\,\mathrm d\gamma + \sum_\gamma\int_\gamma\frac{\alpha\mu}{2|\gamma|}|[u]|^2\mathrm d\gamma. \label{eq:sipdg-stokes-xyz}
\end{align}
In this way, we can use finite element bases that do not meet the continuity requirements of the Stokes equations.


% -------------------------------------
\section{Terrain-following coordinates}

A moving free surface presents an additional challenge because the geometry of the problem is changing in time.
We can solve this problem -- while creating a host of other problems -- by transforming the coordinates.
The idea of terrain-following coordinates is to map the lower surface to 0 and the upper surface to 1.
In so doing, we can work on the Cartesian product of a footprint domain $\Phi$ and the unit interval $[0, 1]$.
We will denote the Cartesian coordinates as $x_1, x_2, x_3$ and terrain-following coordinate as $\xi_1, \xi_2, \xi_3$.
The transformation from terrain-following to Cartesian coordinates is
\begin{equation}
    \left[\begin{matrix}x_1 \\ x_2 \\ x_3\end{matrix}\right] = \left[\begin{matrix}\xi_1 \\ \xi_2 \\ b(\xi_1, \xi_2) + \xi_3h(\xi_1, \xi_2)\end{matrix}\right].
\end{equation}
The complications come in because the surfaces of constant $\xi_3$ are curved.
Taking a derivative with respect to the horizontal coordinates with the elevation $x_3$ fixed is not the same as taking a derivative with $\xi_3$ fixed.
We'll show below how to transform the Stokes equations into this new coordinate system.
In principle, we could compute the metric tensor and then use standard formulas from Riemannian geometry.
Many earth scientists will be unfamiliar with this subject, so we instead derive the transformations explicitly.

The explicit expression for the transformation from terrain-following to Cartesian coordinates will let us also transform vector and tensor fields and their derivatives.
First, we'll introduce a vector field
\begin{equation}
    \sigma = \nabla b + \xi_3\nabla h
\end{equation}
describing the slope in physical space of surfaces of constant $\xi_3$.
The vector field $\sigma$ interpolates between the slope of the bed and the slope of the surface.
In the expression for the slope field, we don't have to specify whether we are taking the gradient in one coordinate system or the other because $b$ and $h$ only depend on the horizontal coordinates.
Then we can concisely express the derivative of the transformation from Cartesian to terrain-following coordinates as
\begin{equation}
    J = \frac{\mathrm dx}{\mathrm d\xi} = \left[\begin{matrix}I & 0 \\ \sigma^* & h\end{matrix}\right].
    \label{eq:dx-dxi}
\end{equation}
We now know how to transform integrals into terrain-following coordinates:
\begin{equation}
    \mathrm dx = |\det J|\mathrm d\xi = h\,\mathrm d\xi.
\end{equation}
As a sanity check, this weighting factor makes the physical units work out correctly.

Now we also have to think about how to transform vectors and tensors between coordinate systems.
The velocities in Cartesian and terrain-following coordinates are not the same.
We can derive the relation between them using the chain rule, and reasoning informally that a velocity is a rate of change of position with time:
\begin{equation}
    u_x = \frac{\mathrm dx}{\mathrm dt} = \frac{\mathrm dx}{\mathrm d\xi}\frac{\mathrm d\xi}{\mathrm dt} = Ju_\xi.
\end{equation}
The body forces transform the same way as well:
\begin{equation}
    f_x = Jf_\xi.
\end{equation}
Intuitively, the gravitational acceleration vector is perpendicular to surfaces of constant depth but is not perpendicular to surfaces of constant relative depth.
We can see from the explicit expression for $J$ in equation \eqref{eq:dx-dxi} that, because of the identity block, the only difference is in the vertical components of the velocity and force.

Next we need to be able to transform derivatives.
To do that, we'll need to know what the derivative of the inverse transformation back from terrain-following to Cartesian coordinates is.
We could compute the inverse transformation and then differentiate it, or we could use the implicit function theorem and the fact that $J$ is a triangular matrix.
These yield the explicit expression
\begin{equation}
    J^{-1} = \frac{\mathrm d\xi}{\mathrm dx} = \left[\begin{matrix}I & 0 \\ -h^{-1}\sigma^* & h^{-1}\end{matrix}\right].
\end{equation}
Given an arbitrary field $\phi$, the chain rule tells us that
\begin{equation}
    \frac{\mathrm d\phi}{\mathrm dx} = \frac{\mathrm d\phi}{\mathrm d\xi}\cdot\frac{\mathrm d\xi}{\mathrm dx} = \frac{\mathrm d\phi}{\mathrm d\xi}\cdot J^{-1}.
\end{equation}
This is just a restatement of the fact that derivatives transform as covectors and not as vectors.
We can then write down an expression for the velocity gradient:
\begin{equation}
    \nabla_xu_x = \nabla_\xi(Ju_\xi)J^{-1}.
\end{equation}
We obtain the strain rate tensor by symmetrizing this expression.

Finally, we need to compute the divergence of a vector field.
In Cartesian coordinates, the divergence operator is the (negative) adjoint of the gradient operator:
\begin{equation}
    \int_\Omega q\nabla_x \cdot v_x\,\mathrm dx = -\int_\Omega\nabla_x q\cdot v_x\,\mathrm dx.
\end{equation}
The general expression for the divergence in other coordinate systems follows by requiring this relationship to remain true.
We know how $\mathrm dx$, $v_x$, and $\nabla_x q$ transform from Cartesian to terrain-following coordinates:
\begin{align}
    \ldots & = -\int_{\Phi \times [0, 1]}\nabla_\xi q \,J^{-1}\cdot Ju_\xi\;h\,\mathrm d\xi \nonumber\\
    & = -\int_{\Phi\times[0, 1]}\nabla_\xi q\cdot h u_\xi\;\mathrm d\xi \nonumber\\
    & = \int_{\Phi\times[0, 1]}q\nabla_\xi\cdot(hu_\xi)\,\mathrm d\xi.
\end{align}
Here we used the fact that the factors of $J^{-1}$ and $J$ cancel out, then applied the usual trick to move the gradient back over as a divegence.
Now we can assert that defining
\begin{equation}
    \nabla_x\cdot u_x = h^{-1}\nabla_\xi\cdot (hu_\xi)
\end{equation}
preserves the relationship between the divergence of a vector field and the gradient of a scalar field.

Now we have all of the ingredients that we need to express the minimization form of the Stokes problem in terrain-following coordinates:
\begin{equation}
    G(u, p) = \int_\Phi\int_0^1\left\{\frac{1}{2}h\tau :\dot\varepsilon - p\nabla\cdot(hu) - hf\cdot J^*J u\right\}\mathrm d\xi
\end{equation}
where the deviatoric stress $\tau$ is again equal to $2\mu\dot\varepsilon$, but the rate of strain tensor is
\begin{equation}
    \dot\varepsilon = \frac{1}{2}\left\{\nabla(Ju)J^{-1} + J^{-*}\nabla (Ju)^*\right\}.
    \label{eq:strain-rate-tfc}
\end{equation}
A reader who knows Riemannian geometry will recognize that the metric tensor $J^*J$ has made its entrance in the inner product between the body force and the velocity.



% -------------------------------------------------------
\section{Discretization in terrain-following coordinates}

Suppose that we discretize the bed elevation $b$ and thickness $h$ using a standard continuous finite element basis.
The derivative $J$ of the coordinate transformation involves gradients of $b$ and $h$.
That means that $J$ is discontinuous across cell boundaries.
But if we examine expression \eqref{eq:strain-rate-tfc}, $J$ appears under a gradient and we are faced with the problem of how to differentiate a discontinuous quantity.
This lack of continuity is the major problem we need to overcome.
We can solve that problem with the interior penalty DG method.
At a practical level, this amounts to substituting the transformed expressions for $u$ and its derivatives in terrain-following coordinates into the functional in equation \eqref{eq:sipdg-stokes-xyz}.


% ----------------------------------
\section{Testing the implementation}

There are a lot of moving parts in the derivations above.
We have designed a battery of tests to catch potential failures of either the math or the code.
Our first level of testing is to isolate whether our implementation of the interior penalty DG method is correct.
Implementing this discretization is notoriously error-prone.
If our solution in terrain-following coordinates is incorrect, it's not immediately apparent whether that's because of a failure in the use of IPDG or in its application in a different coordinate system.
Testing IPDG alone first can help to pinpoint the source of later failures.
Once we have some assurances that the IPDG discretization works as expected, we will then test the discretization in terrain-following coordinates.
We will carry out each of these steps on a sequence of meshes and check for convergence at the expected rates.

First, we solve the Stokes equations in Cartesian coordinates with variable surface and bed topography using a conforming finite element basis.
Then we will use discontinuous basis functions and the symmetric interior-penalty method.
If the norm difference between the two solutions does not go to zero as the mesh is refined, then the most likely explanation is that either the math or the implementation of the interior penalty method is incorrect.

The real test comes when we transform the problem to terrain-following coordinates.
If we map the solution in terrain-following coordinates back to Cartesian coordinates and obtain a vector field that differs appreciably from what we obtained in the previous steps, then we will know that either the math or the implementation is wrong.


\pagebreak

\bibliographystyle{plainnat}
\bibliography{stokes-tfc.bib}

\end{document}
